\documentclass[a4paper,11pt]{article}

% 1. 中文支持 (必须使用 XeLaTeX 编译)
% quiet: 抑制 fontspec 对 TeX Live 自带 Fandol 字体的无害 "Script CJK" 警告
\PassOptionsToPackage{quiet}{fontspec}
% xcolor 选项需在 tikz/ctex 载入 xcolor 之前传入,否则 option clash
\PassOptionsToPackage{usenames,dvipsnames}{xcolor}
\usepackage[UTF8]{ctex}

% 2. 绘图包 (实现蓝色梯形标题)
\usepackage{tikz}
\usetikzlibrary{shapes, backgrounds}

% 3. 布局与格式包
\usepackage{latexsym}
\usepackage[empty]{fullpage}
\usepackage[explicit]{titlesec}
\usepackage{marvosym}
\usepackage{verbatim}
\usepackage{enumitem}
\usepackage[hidelinks]{hyperref}
\usepackage{fancyhdr}
\usepackage{tabularx}
\usepackage{xcolor}

%-------------------------------------------
% 颜色定义
\definecolor{MyDarkBlue}{RGB}{0, 50, 100}

% 页边距设置 (优化后让页面更饱满)
\pagestyle{fancy}
\fancyhf{}
\fancyfoot{}
\renewcommand{\headrulewidth}{0pt}
\renewcommand{\footrulewidth}{0pt}

\addtolength{\oddsidemargin}{-0.6in}
\addtolength{\evensidemargin}{-0.6in}
\addtolength{\textwidth}{1.2in}
\addtolength{\topmargin}{-.7in}
\addtolength{\textheight}{1.4in}

\urlstyle{same}
\raggedbottom
\raggedright
\setlength{\tabcolsep}{0in}

% --- 自定义 TikZ 标题样式 ---
\titleformat{\section}
  {}
  {}
  {0pt}
  {%
    \noindent
    \begin{tikzpicture}[baseline=(t.base)]
    \vspace{-2pt}
      \node[inner xsep=10pt, inner ysep=3pt, text=white, font=\large\bfseries, align=left, anchor=base west] (t) at (0,0) {#1};
      \begin{scope}[on background layer]
         \fill[MyDarkBlue] (t.north west) -- (t.north east) -- ([xshift=8pt]t.south east) -- (t.south west) -- cycle;
         \draw[MyDarkBlue, line width=2pt] ([xshift=-68.30pt]t.south east) -- (\linewidth, 0 |- t.south east);
      \end{scope}
    \end{tikzpicture}%
  }
  []
\titlespacing*{\section}{0pt}{12pt}{6pt}

%-------------------------------------------
% 自定义命令
\newcommand{\resumeItem}[1]{
  \item\small{{#1 \vspace{-1pt}}}
}

\newcommand{\resumeEducationHeading}[4]{
  \vspace{-2pt}\item
    \begin{tabular*}{0.97\textwidth}[t]{l@{\extracolsep{\fill}}r}
      \textbf{#1} & #2 \\
      \textit{\small #3} & \textit{\small #4} \\
    \end{tabular*}\vspace{-5pt}
}

\newcommand{\resumeResearchHeading}[2]{
  \vspace{-2pt}\item
    \begin{tabular*}{0.97\textwidth}[t]{l@{\extracolsep{\fill}}r}
      \textbf{#1} & \textbf{#2} \\
    \end{tabular*}\vspace{-7pt}
}

\newcommand{\resumeProjectHeading}[2]{
    \item
    \begin{tabular*}{0.97\textwidth}{l@{\extracolsep{\fill}}r}
      \small#1 & #2 \\
    \end{tabular*}\vspace{-7pt}
}

\newcommand{\resumeSubHeadingListStart}{\begin{itemize}[leftmargin=0.15in, label={}]}
\newcommand{\resumeSubHeadingListEnd}{\end{itemize}}
\newcommand{\resumeItemListStart}{\begin{itemize}[leftmargin=*, label={\small\textbullet}]}
\newcommand{\resumeItemListEnd}{\end{itemize}\vspace{-2pt}}

%-------------------------------------------
%%%%%%  简历正文  %%%%%%%%%%%%%%%%%%%%%%%%%%%%
\begin{document}

%---------- 个人信息 ----------
\begin{center}
    \textbf{\Huge {\ziju{0.2} 薛景秦}} \\ \vspace{8pt}
    \small
    \href{mailto:xuejingqin@ihep.ac.cn}{\underline{xuejingqin@ihep.ac.cn}} $|$
    +86 151-1030-9796 $|$
    中国科学院大学 $\cdot$ 高能物理研究所 $|$
    
\end{center}

%----------- 教育背景 -----------
\section{教育背景}
  \resumeSubHeadingListStart
    \resumeEducationHeading
      {中国科学院大学 \ (高能物理研究所)}{2023.09 -- 2023.06}
      {硕士 \ \textbar \ 粒子物理与原子核物理}{北京}
      \resumeItemListStart
        \resumeItem{\textbf{GPA: 3.2 / 4.0} \ \ $|\ \ $ \textbf{主要荣誉：}三好学生、国家奖学金、所长奖学金}
        \resumeItem{\textbf{一作论文：}发表于 \textit{Eur. Phys. J. C} 85 (2025) 1420（IF 4.8，中科院二区，作者排名第一）}
        \resumeItem{\textbf{英语水平：}CET-4 / CET-6}
      \resumeItemListEnd
    \resumeEducationHeading
      {中北大学 \ (半导体与物理学院)}{2019.09 -- 2023.06}
      {本科 \ \textbar \ 应用物理学}{太原，山西}
      \resumeItemListStart
        \resumeItem{\textbf{GPA: 4.30 / 5.0} \ \ $|\ \ $ \textbf{专业排名: 1 / 100 (Top 1\%)} \ \ $|\ \ $ \textbf{保送研究生}}
        \resumeItem{校级“优秀毕业生”、校长奖学金提名奖（全院仅 20 名）、国家奖学金}
        \resumeItem{连续两届全国大学生数学竞赛省一等奖、全国大学生数学建模竞赛省二等奖、美国大学生数学建模竞赛 H 奖}
      \resumeItemListEnd
  \resumeSubHeadingListEnd

%----------- 个人技能 -----------
\section{个人技能}
 \begin{itemize}[leftmargin=0.15in, label={}]
    \small{\item{
     \textbf{编程语言}{: \textbf{C/C++}（熟练，掌握内存管理与底层优化）、Python、Pytorch、Shell 脚本；LaTeX} \\
     \textbf{高性能计算 (HPC)}{: \textbf{PyTorch (autograd)} 张量化建模与批处理、算法\textbf{张量化表示}重构、多核 / 集群并行、分布式作业提交} \\
     \textbf{性能调优}{: Linux \texttt{perf}、Hotspot 性能剖析；热点函数级底层优化} \\
     \textbf{工程工具}{: Linux 开发环境、Git 协同开发、CMake / Makefile 构建系统、Claude Code 等 Agentic 编程工具} \\
     \textbf{数据与统计}{: 大规模蒙特卡洛模拟、最大似然估计与 $\chi^2$ 拟合、假设检验与置信区间构造 (Feldman--Cousins)、覆盖率检验与系统误差建模}
    }}
 \end{itemize}

%----------- 科研与项目经历 -----------
\section{科研与项目经历}
  \resumeSubHeadingListStart

    \resumeResearchHeading
      {高性能前向折叠能谱拟合框架的开发与加速实现}{2024.09 -- 2025.05}
      \resumeItemListStart
        \resumeItem{\textbf{面向对象的 PyTorch 拟合框架架构设计与性能优化} \hfill \textbf{\color{MyDarkBlue}[独立设计 / 核心开发 / 第一作者]}}
        \resumeItem{从零搭建框架架构，运用\textbf{面向对象编程}思想将“反应堆能谱 $\to$ IBD 微分截面 $\to$ 沉积能量谱 $\to$ 探测器响应 $\to$ 重建能谱”的多层处理链路抽象为\textbf{可复用、可扩展的模块化类层次}：各物理算子封装为独立组件，经统一接口级联为\textbf{四级张量算子流水线}，显著降低耦合度并便于扩展与测试。}
        \resumeItem{基于 \textbf{PyTorch autograd} 将整条链路重构为端到端可微的\textbf{张量计算图}，把数十项系统误差统一编码为模型参数，以\textbf{梯度下降（gradient descent）}做 $\chi^2$ 极小化，替代传统数值求导，大幅提升收敛速度与数值稳定性。}
        \resumeItem{基于 Profiling 自顶向下定位性能热点（信号谱计算占 98.6\%），实施三项优化：\textbf{分层缓存}（静态量预计算 + 函数级 memoization，约 3$\times$）、\textbf{带状稀疏矩阵}（CSR 存储，矩阵-向量乘 $O(N^2)\to O(N\cdot w)$，约 1.5$\times$）、\textbf{分块张量化构造}（块宽 512，临时张量驻留 L2 缓存，约 1.5$\times$）。}
        \resumeItem{单次拟合耗时：JUNO only \textbf{110\,s $\to$ 15.2\,s}、JUNO + DYB 联合 \textbf{50\,min $\to$ 5\,min}，整体加速超一个数量级；双精度基准下相对误差 $< 10^{-6}$，5000 组 Toy-MC 收敛成功率达 100\%。}
        \resumeItem{框架支撑 \textbf{500 万次拟合}的大规模计算，成果以\textbf{第一作者}发表于 \textbf{\textit{Eur. Phys. J. C} 85 (2025) 1420}（IF 4.8，中科院二区），被 JUNO 合作组采纳为\textbf{基础分析工具}，并在第 25 届 JUNO 国际合作会议作主题报告。}
      \resumeItemListEnd

    \resumeResearchHeading
      {大规模 MC 模拟与统计方法研究}{2025.05 -- 2025.11}
      \resumeItemListStart
        \resumeItem{\textbf{基于面向对象 Python 的大规模 Toy-MC 模拟与统计分析} \hfill \textbf{\color{MyDarkBlue}[核心开发]}}
        \resumeItem{基于\textbf{面向对象的 Python} 构建可配置、可复用的自动化 Toy-MC 流水线，支持万级伪实验、多种分 bin 方案一键切换；以分位数（第 16/84 百分位）非参数化定义偏差指标，规避高斯近似假设。}
        \resumeItem{\textbf{（1）能谱分 bin 策略：}将重建能谱 bin 宽由 20\,keV 调整至 \textbf{100\,keV}，关键参数统计偏差由 0.8$\sigma$ 降至 $<$ \textbf{0.2$\sigma$}，模型依赖系统误差最小（$<$ \textbf{0.3\%}）—— 被 JUNO 首篇物理论文采纳为分 bin 策略（Nature, arXiv:2511.14593）。}
        \resumeItem{\textbf{（2）Feldman--Cousins 方法工程化：}将 FC 构造分解为 50 网格点 $\times$ 5{,}000 伪实验 $\times$ 20 多起始点 $=$ \textbf{500 万次独立拟合}，在 \textbf{1{,}000+ 核 HPC 集群}上并行编排；针对似然面多极小值采用多起始点全局优化。}
        \resumeItem{由经验 $\Delta\chi^2$ 分布提取临界值 \textbf{$\Delta\chi^2 \approx 1.57$}（显著高于渐近近似 1.0），定量证实 Wilks 定理在低统计量下失效；相关统计策略建议被 JUNO 首篇物理论文采纳（Nature, arXiv:2511.14593）。}
      \resumeItemListEnd

    \resumeResearchHeading
      {OMILREC 事例重建算法的加速实现}{2025.06 -- 2025.09}
      \resumeItemListStart
        \resumeItem{\textbf{基于面向对象 C++ 的事例重建算法性能剖析与底层优化} \hfill \textbf{\color{MyDarkBlue}[核心开发]}}
        \resumeItem{在\textbf{面向对象 C++} 重建框架中，基于 Linux \texttt{perf} + Hotspot 性能剖析自顶向下定位关键热点函数，针对性重构核心计算模块并实施底层优化。}
        \resumeItem{引入\textbf{稀疏化计算}，按事例信息筛选有效 PMT、跳过无贡献 PMT，避免对全部 PMT 的无效遍历。}
        \resumeItem{通过\textbf{多线程并行}与几何计算\textbf{向量化（SIMD）}进一步压缩单次重建耗时。}
        \resumeItem{将与拟合参数无关的计算\textbf{移出 Minuit 似然循环}，预计算后复用，消除迭代中的重复开销；最终\textbf{速度提升 3 倍}，并在第 26 届 JUNO 国际合作会议作主题报告。}
      \resumeItemListEnd

  \resumeSubHeadingListEnd

%----------- 学术成果与会议报告 -----------
\section{学术成果与会议报告}
  \resumeItemListStart
    \resumeItem{\textbf{[期刊一作]} \textbf{J.~Xue}, X.~Ding, et al., ``High-Performance Statistical Methods for Reactor Neutrino Oscillations,'' \textit{Eur. Phys. J. C} \textbf{85} (2025) 1420.}
    \resumeItem{\textbf{[会议一作]} \textbf{J.~Xue}, ``The accelerated reactor neutrino fitter and oscillation analysis of JUNO,'' \textit{PoS}.}
    \resumeItem{\textbf{[合作组文章]} A.~Abusleme, T.~Adam, \ldots, \textbf{J.~Xue}, et al., ``First measurement of reactor neutrino oscillations at JUNO,'' \textit{Nature} (accepted, 2026); arXiv:2511.14593.}
    \resumeItem{\textbf{[合作组文章]} A.~Abusleme, T.~Adam, \ldots, \textbf{J.~Xue}, et al., ``Initial performance results of the JUNO detector,'' \textit{Chinese Physics C} (2026).}
    \resumeItem{\textbf{[主题报告]} 第 25 届江门中微子国际合作会议：\textit{The accelerated reactor fitter}。}
    \resumeItem{\textbf{[主题报告]} 第 26 届江门中微子国际合作会议：\textit{Reducing Bias in Low Exposure: Poisson Method and Binning Strategies}。}
    \resumeItem{\textbf{[主题报告]} 第 26 届江门中微子国际合作会议：\textit{Speedup of OMILREC}。}
    \resumeItem{\textbf{[Poster]} TAUP 2025 国际会议：\textit{The accelerated reactor fitter and Oscillation analysis for JUNO}。}
  \resumeItemListEnd

%----------- 荣誉奖项 -----------
\section{荣誉奖项}
  \resumeSubHeadingListStart
    \resumeProjectHeading
      {\textbf{国家奖学金}（教育部，研究生 \& 本科共获两次，Top 0.2\% 顶尖荣誉）}{研究生 / 本科}
    \resumeProjectHeading
      {\textbf{三好学生} \ \& \ \textbf{所长奖学金}（中国科学院大学 / 高能物理研究所）}{研究生}
    \resumeProjectHeading
      {\textbf{校级“优秀毕业生”} \ \& \ \textbf{校长奖学金提名奖}（全院仅 20 名，Top 0.2\% 顶尖荣誉）}{本科}
    \resumeProjectHeading
      {\textbf{全国大学生数学竞赛} 省一等奖（连续两届）}{本科}
    \resumeProjectHeading
      {\textbf{美国大学生数学建模竞赛 (MCM/ICM)} H 奖 \ \& \ \textbf{全国大学生数学建模竞赛} 省二等奖}{本科}
  \resumeSubHeadingListEnd

\end{document}
